% !TEX root = ../main.tex

\section{Twierdzenie Clausiusa}
Rozważamy dowolny cykl zamknięty \(X\), w którym temperatura jest ograniczona od góry temperaturą \(T_0\), to znaczy, że dla każdego punktu cyklu \(T \leq T_0\). Pomiędzy zbiornikiem o temperaturze \(T_0\) a gazem roboczym silnika \(X\) dla każdego punktu jego cyklu umieszczamy silnik Carnota. Dodatni przepływ ciepła wybieramy jako: ciepło pobrane ze zbiornika o temperaturze \(T_0\) przez silnik Carnota \(\idif{Q_0}\); ciepło oddane przez silnik Carnota do silnika \(X\). Podczas pracy silników Carnota gaz w silniku \(X\) ma stałą temperaturę, ponieważ przez każdy z nich dostarczane jest jedynie ciepło elementarne \(\idif{Q}\). Z porównania równań \eqref{eq:sprawnosc_carnot_Q} i \eqref{eq:sprawnosc_carnot} otrzymujemy dla ciepeł elementarnych:
\begin{equation*}
	\frac{\idif{Q_0}}{T_0} = \frac{\idif{Q}}{T}.
\end{equation*}
Całkowite ciepło dostarczone ze zbiornika o temperaturze \(T_0\) jest równe całce po przebiegu silnika \(X\):
\begin{equation*}
	Q_0 = \oint \idif{Q_0} = T_0 \oint \frac{\idif{Q}}{T}.
\end{equation*}
Jeśli \(Q_0 > 0\) znaczyłoby to, że całkowite ciepło zostało przekształcone na pracę, która jest sumą pracy silników Carnota i \(X\). Stąd wymagane jest, aby:
\begin{equation}\label{eq:twierdzenie_clausius}
	\oint \frac{\idif{Q}}{T} \leq 0.
\end{equation}
Jest to twierdzenie Clausiusa.

Dla cyklu \(X\) odwracalnego zmieniamy kierunki wszystkich przepływów, stąd:
\begin{equation*}
	\oint \frac{\idif{Q^R}}{T} \geq 0.
\end{equation*}
Zarówno ta nierówność jak i \eqref{eq:twierdzenie_clausius} muszą być spełnione, musi więc zachodzić:
\begin{equation*}
	\oint \frac{\idif{Q^R}}{T} = 0.
\end{equation*}
Znaczy to, że wielkość \(\frac{\idif{Q^R}}{T}\) jest różniczką zupełną:
\begin{equation}\label{eq:def_entropia}
	\d{S} \overset{def}{=} \frac{\idif{Q^R}}{T}.
\end{equation}
W układzie izolowanym nie ma przepływów ciepła i entropia spełnia:
\begin{equation}\label{eq:2zt_dS}
	\d{S} \geq 0.
\end{equation}
Całkując powyższą definicję po drodze odwracalnej mamy:
\begin{equation*}
	S\pab{B} - S\pab{A} = \sideset{_R}{_A^B}{\int} \frac{\idif{Q}}{T}.
\end{equation*}

W przypadku ogólnym możemy rozpisać całkę \eqref{eq:twierdzenie_clausius} na dwie części:
\begin{gather*}
	0 \geq \oint \frac{\idif{Q}}{T} = 
	\sideset{_N}{_A^B}{\int} \frac{\idif{Q}}{T} + \sideset{_R}{_B^A}{\int} \frac{\idif{Q}}{T} = 
	\sideset{_N}{_A^B}{\int} \frac{\idif{Q}}{T} - \sideset{_R}{_A^B}{\int} \frac{\idif{Q}}{T} \\
	\sideset{_N}{_A^B}{\int} \frac{\idif{Q}}{T} \leq \sideset{_R}{_A^B}{\int} \frac{\idif{Q}}{T} = S\pab{B} - S\pab{A},
\end{gather*}
gdzie indeks dolny \(R\) i \(N\) przy całce oznaczają odpowiednio przemianę odwracalną oraz dowolną (w tym nieodwracalną). Zmiana entropii związana z przepływem ciepła jest równa:
\begin{equation}\label{eq:def_deS}
	\d_e{S} \overset{def}{=} \frac{\idif{Q}}{T}.
\end{equation}
Wykorzystując tą wielkość otrzymujemy:
\begin{gather*}
	\d{S} = \d_i{S} + \d_e{S} \geq \d_e{S}, \\
	\d_i{S} \geq 0.
\end{gather*}
Jest to przypadek ogólny i razem z \eqref{eq:2zt_dS} stanowi matematyczne sformuowanie Drugiej Zasady Termodynamiki.